% !TEX program = pdflatex
% arara: pdflatex: { shell: false }
% arara: pdflatex: { shell: false }
% arara: latexmk: { options: ["-c"] }
% -------------------------------
\input{inc/_m.latex}
% \ShowTOCfalse
% -------------------------------
\nc{\Classe}{2nde}
\nc{\Titre}{Droites du plan}
% -------------------------------

\begin{document}
\Entete
\section{Vecteur directeur et équation cartésienne d'une droite}
\subsection{Vecteur directeur}

\begin{defin}{Vecteur directeur}
	\begin{wrapfigure}{r}{6cm}
		\begin{center}
			\begin{tikzpicture}[scale=1.1]
				\draw[coulAccent, thick] (-1.5,-1) -- (3.5,2.333) node[right] {$d$};
				% petit vecteur
				\draw[-{Stealth[length=10pt, width=4pt]}, coulAlerte, ultra thick]
				(-1,-0.667) -- (0.5,0.333) node[midway, above left] {$\vec{u}$};
				% vecteur moyen
				\draw[-{Stealth[length=10pt, width=4pt]}, coulVert, ultra thick]
				(0.5,1.333) -- (2.5,2.667) node[midway, above left] {$\vec{v}$};
				% grand vecteur dans l'autre sens
				\draw[-{Stealth[length=10pt, width=4pt]}, coulOrange, ultra thick]
				(3.5,1.667-1) -- (2,0.667-1) node[midway, below right] {$\vec{w}$};
			\end{tikzpicture}
		\end{center}
	\end{wrapfigure}
	Soit $d$ une droite du plan.

	On appelle \textbf{vecteur directeur} de $d$ tout vecteur non nul $\vec{u}$ qui possède la même direction que la droite $d$

	\vspace*{3.5cm}
\end{defin}

\begin{methode}{Déterminer graphiquement un vecteur directeur d'une droite}
	Donner des vecteurs directeurs des droites $d_1$, $d_2$, $d_3$ et $d_4$

	\begin{center}
		\begin{tikzpicture}[scale=0.8]
			\draw[very thin, coulBordure!50] (-0.9,-0.9) grid (6.9,6.9);
			\draw[-{Stealth}, coulPrimaire] (-0.5,0) -- (7,0) node[right] {$x$};
			\draw[-{Stealth}, coulPrimaire] (0,-0.5) -- (0,7) node[above] {$y$};
			% d1 : pente 2
			\draw[coulAccent, ultra thick] (-1,-1) -- (3.5,7) node[below left, xshift=-3mm] {$d_1$};
			% d2 : horizontale
			\draw[coulVert, ultra thick] (-1,3) -- (7,3) node[above left] {$d_2$};
			% d3 : pente -1
			\draw[coulOrange, ultra thick] (-1,7) -- (7,-1) node[above right] at (6,0) {$d_3$};
			% d4 : verticale
			\draw[mulberry, ultra thick] (4,-1) -- (4,7) node[right, yshift=-3mm] {$d_4$};
			% ticks
			\foreach \x in {1,...,6} \node[below, font=\tiny, coulPrimaire] at (\x,0) {\x};
			\foreach \y in {1,...,6} \node[left, font=\tiny, coulPrimaire] at (0,\y) {\y};
		\end{tikzpicture}
	\end{center}
\end{methode}
\newpage
\begin{methode}{Déterminer graphiquement un vecteur directeur d'une droite}
	\begin{correction}
		\bi
		\item Pour $d_1$ : $\vec{a}\coord{1}{2}$ convient.
		On peut aussi prendre $\vec{b}\coord{2}{4}$ ou $\vec{c}\coord{-1}{-2}$
		\item Pour $d_2$ : $\vec{u}\coord{6}{0}$ convient (droite horizontale).
		\item Pour $d_3$ : $\vec{v}\coord{1}{-1}$ convient.
		\item Pour $d_4$ : $\vec{w}\coord{0}{2}$ convient (droite verticale).
		\ei
	\end{correction}
\end{methode}

\subsection{Équation cartésienne d'une droite}

\begin{defin}{Équation cartésienne}
	Toute droite admet une équation de la forme : $$\boxed{ax + by + c = 0}$$ \ldots avec $(a\,;b) \neq (0\,;0)$

	\ms

	Cette équation est appelée \textbf{équation cartésienne} de la droite.
\end{defin}

\begin{prop}{Vecteur directeur et équation cartésienne}
	Le vecteur $\vec{u}\coord{-b}{a}$ est un vecteur directeur de la droite d'équation cartésienne $ax + by + c = 0$
\end{prop}

\begin{demo}{au programme}
	Soit $A\coordl{x_0}{y_0}$ un point de la droite $d$ et $\vec{u}\coord{\alpha}{\beta}$ un vecteur directeur de $d$

	Un point $M\coordl{x}{y}$ appartient à $d$ si et seulement si $\vec{AM}\coord{x-x_0}{y-y_0}$ et $\vec{u}\coord{\alpha}{\beta}$ sont colinéaires, soit $\det(\vec{AM}\,;\vec{u}) = 0$, c'est-à-dire :

	\[
		\determinant{x-x_0 & \alpha \\ y-y_0 & \beta} = 0
		\Rarr \beta(x-x_0) - \alpha(y-y_0) = 0.
	\]

	En développant :

	$$
		\begin{array}{crcl}
			~    & \beta x - \beta x_0 - \alpha y + \alpha y_0   & = & 0 \\
			\eqv & \beta x - \alpha y + (\alpha y_0 - \beta x_0) & = & 0 \\
			\eqv & ax - by + (\alpha y_0 - \beta x_0)            & = & 0 \\
		\end{array}
	$$

	Cette équation s'écrit $ax + by + c = 0$ avec :
	\bi
	\item $a = \beta$
	\item $b = -\alpha$
	\item $c = \alpha y_0 - \beta x_0$
	\ei

	Les coordonnées de $\vec{u}$ sont donc $\coord{\alpha}{\beta} = \coord{-b}{a}$
\end{demo}

\newpage

\begin{ex}{Vecteur directeur depuis l'équation cartésienne}
	Un vecteur directeur de la droite d'équation $4x - 5y - 1 = 0$ est $\vec{u}\coord{5}{4}$

	En effet, $a = 4$ et $b = -5$ donc $\vec{u}\coord{-b}{a} = \coord{5}{4}$
\end{ex}

\begin{methode}{Déterminer une équation cartésienne depuis un point et un vecteur directeur}
	\begin{enumerate}
		\item Déterminer une équation cartésienne de $d$ passant par $A\coordl{3}{1}$ et de vecteur directeur $\vec{u}\coord{-1}{5}$
		\item Déterminer une équation cartésienne de $d'$ passant par $B\coordl{5}{3}$ et $C\coordl{1}{-3}$
	\end{enumerate}

	\begin{correction}

		\textbf{1)} $d$ admet une équation de la forme $ax + by + c = 0$

		$\vec{u}\coord{-1}{5} = \coord{-b}{a}$ donc $a = 5$ et $b = 1$

		L'équation est de la forme $5x + y + c = 0$

		On substitue les coordonnées de $A$ :
		\[
			\begin{array}{rccccccc}
				     & 5 \times x_A & + & y_A & + & c & = & 0   \\
				\eqv & 5 \times 3   & + & 1   & + & c & = & 0   \\
				\eqv &              &   & 16  & + & c & = & 0   \\
				\eqv &              &   &     &   & c & = & -16
			\end{array}
		\]
		\ms
		Une équation de $d$ est $5x + y - 16 = 0$

		\ms

		\textbf{2)} $\vec{BC}\coord{1-5}{-3-3} = \coord{-4}{-6} = \coord{-b}{a}$ donc $a = -6$ et $b = 4$

		L'équation est de la forme $-6x + 4y + c = 0$

		On substitue les coordonnées de $B$ :
		\[
			\begin{array}{rccccccc}
				     & -6 \times x_B & + & 4 \times y_B & + & c & = & 0  \\
				\eqv & -6 \times 5   & + & 4 \times 3   & + & c & = & 0  \\
				\eqv &               &   & -18          & + & c & = & 0  \\
				\eqv &               &   &              &   & c & = & 18
			\end{array}
		\]

		Une équation cartésienne de $d'$ est $-6x + 4y + 18 = 0$, soit $-3x + 2y + 9 = 0$
	\end{correction}
\end{methode}

\begin{methode}{Tracer une droite depuis son équation cartésienne}
	Tracer la droite $d$ d'équation cartésienne $3x + 2y - 5 = 0$

	\begin{correction}

		On cherche un point et un vecteur directeur.

		\bi
		\item Pour $x = 0$ :

		$$
			\begin{array}{rl}
				     & 3\times 0 + 2y - 5 = 0   \\
				\eqv & 2y - 5 = 0               \\
				\eqv & y = \dfrac{5}{2} = 2{,}5
			\end{array}
		$$

		Le point $A\coordl{0}{2{,}5}$ appartient à $d$

		\item $a = 3$, $b = 2$ donc $\vec{u}\coord{-2}{3}$ est un vecteur directeur.

		\ei
	\end{correction}
\end{methode}
\newpage

\begin{methode}{}
	\begin{correction}
		\begin{center}
			\begin{tikzpicture}[scale=0.8]
				\draw[very thin, coulBordure!50] (-3.5,-2.5) grid (5.5,6.5);
				\draw[-{Stealth}, coulPrimaire] (-3,0) -- (5,0) node[right] {$x$};
				\draw[-{Stealth}, coulPrimaire] (0,-2) -- (0,6) node[above] {$y$};
				\foreach \x in {-3,...,4} \node[below, font=\tiny] at (\x,0) {\x};
				\foreach \y in {-2,...,5} \node[left, font=\tiny] at (0,\y) {\y};
				% droite 3x+2y-5=0, y = -3/2 x + 5/2
				\draw[coulAccent, ultra thick, domain=-2.5:3] plot(\x, {-1.5*\x + 2.5})
				node[right] {$d$};
				% point A
				\filldraw[coulAlerte] (0,2.5) circle (0.08)
				node[right, font=\small] {$A(0\,;\,2{,}5)$};
				% point B 
				\filldraw[coulAlerte] (-2,5.5) circle (0.08);
				% vecteur directeur
				\draw[-{Stealth[length=11pt]}, coulVert, ultra thick] (4,1) -- +(-2,3) node[midway, right, font=\small] {$\vec{u}$};
				\draw[-{Stealth[length=10pt]}, coulVert, ultra thick] (0,2.5) -- +(-2,3);
			\end{tikzpicture}
		\end{center}
	\end{correction}
\end{methode}

\subsection{Position relative de deux droites}

\begin{prop}{Droites parallèles et vecteurs directeurs colinéaires}
	Deux droites sont parallèles si et seulement si elles ont des vecteurs
	directeurs colinéaires.
\end{prop}

\begin{methode}{Déterminer la position relative de deux droites}
	Démontrer que les droites $d_1 : 6x - 10y - 5 = 0$ et
	$d_2 : -9x + 15y = 0$ sont parallèles.

	\begin{correction}

		$\vec{u}\coord{10}{6}$ est un vecteur directeur de $d_1$ et
		$\vec{v}\coord{-15}{-9}$ est un vecteur directeur de $d_2$

		\[
			\det(\vec{u}\,;\vec{v}) = \determinant{10 & -15 \\ 6 & -9}
			= 10 \times (-9) - 6 \times (-15) = -90 + 90 = 0.
		\]

		Donc $\vec{u}$ et $\vec{v}$ sont colinéaires, et les droites $d_1$ et
		$d_2$ sont parallèles.
	\end{correction}
\end{methode}

% ══════════════════════════════════════════════
\section{Équation réduite et pente d'une droite}
% ══════════════════════════════════════════════

\subsection{Équation réduite}

\begin{prop}{Forme de l'équation réduite}
	Soit une droite $d$
	\bi
	\item Si $d$ est \textbf{parallèle à l'axe des ordonnées}, son équation
	est de la forme $x = n$
	\item Sinon, son équation est de la forme $y = mx + p$, appelée
	\textbf{équation réduite} de $d$
	\bi
	\item $m$ est la \textbf{pente} (ou \textbf{coefficient directeur}) de $d$
	\item $p$ est l'\textbf{ordonnée à l'origine} de $d$
	\ei
	\ei
\end{prop}

\newpage
\begin{demo}{}
	\bi
	\item Si $b \neq 0$, l'équation $ax + by + c = 0$ donne
	$y = -\dfrac{a}{b}x - \dfrac{c}{b}$, avec $m = -\dfrac{a}{b}$ et
	$p = -\dfrac{c}{b}$
	\item Si $b = 0$, l'équation donne $x = -\dfrac{c}{a}$ : droite parallèle
	à l'axe des ordonnées.
	\ei
\end{demo}

\begin{methode}{Passer d'une équation cartésienne à l'équation réduite}
	\begin{enumerate}
		\item Soit $d : 6x + 3y - 5 = 0$ Déterminer l'équation réduite de $d$
		\item Soit $d' : y = 6x - 5$ Déterminer une équation cartésienne de $d'$
	\end{enumerate}

	\begin{correction}

		\textbf{1)} On isole $y$ :
		\eqn{
			6x + 3y - 5 &= 0 \\
			3y &= -6x + 5 \\
			y &= -2x + \frac{5}{3}
		}
		Équation réduite de $d$ : $y = -2x + \dfrac{5}{3}$

		\ms

		\textbf{2)} On ramène tout à gauche :
		\[
			y = 6x - 5 \eqv -6x + y + 5 = 0.
		\]
		Équation cartésienne de $d'$ : $-6x + y + 5 = 0$
	\end{correction}
\end{methode}

\begin{exercice}{Pente et ordonnée à l'origine}
	Donner la pente et l'ordonnée à l'origine de chacune des droites :
	\begin{enumerate}
		\item $y = -2x + 3$
		\item $y = 5$
		\item $4x + 2y = 1$
	\end{enumerate}

	\begin{correction}
		\begin{enumerate}
			\item Pente : $-2$ \quad Ordonnée à l'origine : $3$
			\item Pente : $0$ \quad Ordonnée à l'origine : $5$
			\item On réécrit :

			      $$\begin{array}{rrcl}
					           & 4x + 2y & = & 1                  \\
					      \eqv & 2y      & = & 1-4x               \\
					      \eqv & y       & = & -2x + \tfrac{1}{2}
				      \end{array}$$

			      \bi\item Pente : $-2$\item Ordonnée à l'origine : $\dfrac{1}{2}$\ei
		\end{enumerate}
	\end{correction}
\end{exercice}

\newpage

\begin{methode}{Tracer une droite depuis son équation réduite}
	Tracer les droites :
	\bi
	\item $d_1 : y = 2x + 3$
	\item  $d_2 : y = 4$
	\item $d_3 : x = 3$
	\ei

	\begin{correction}
		\bi
		\item $d_1$ passe par $(0\,;3)$ et $(2\,;7)$ car :
		\bi
		\item pour $x=0$, on a : $y = 2\times 0+3=3$
		\item pour $x=2$, on a : $y = 2\times 2+3=7$
		\ei
		\item $d_2$ est la droite horizontale d'ordonnée $4$
		\item $d_3$ est la droite verticale d'abscisse $3$
		\ei
		\begin{center}
			\begin{tikzpicture}[scale=0.8]
				\draw[very thin, coulBordure!50] (-0.5,-0.5) grid (5.5,8.5);
				\draw[-{Stealth}, coulPrimaire] (-0.5,0) -- (5.5,0) node[right] {$x$};
				\draw[-{Stealth}, coulPrimaire] (0,-0.5) -- (0,8.5) node[above] {$y$};
				\foreach \x in {1,...,5} \node[below, font=\tiny] at (\x,0) {\x};
				\foreach \y in {1,...,8} \node[left, font=\tiny] at (0,\y) {\y};
				% d1 : y = 2x+3
				\draw[coulAccent, ultra thick, domain=-0.5:3]
				plot(\x, {2*\x+3}) node[left] {$d_1$};
				\filldraw[coulAccent] (0,3) circle (0.1);
				\filldraw[coulAccent] (2,7) circle (0.1);
				% d2 : y = 4
				\draw[coulVert, ultra thick] (-0.5,4) -- (5.5,4) node[right] {$d_2$};
				% d3 : x = 3
				\draw[coulOrange, ultra thick] (3,-0.5) -- (3,8.5) node[right] {$d_3$};
			\end{tikzpicture}
		\end{center}
	\end{correction}
\end{methode}

\begin{methode}{Vérifier si un point appartient à une droite}
	Les points $A\coordl{6}{39}$ et $B\coordl{346}{2420}$ appartiennent-ils à la droite $d : y = 7x - 3$ ?

	\begin{correction}
		\bi
		\item $7 \times 6 - 3 = 42 - 3 = 39$ \\ $\rarr$ les coordonnées de $A$ vérifient l'équation $\quad\rarr\quad$ \textbf{$A$ appartient à $d$}
		\item $7 \times 346 - 3 = 2422 - 3 = 2419 \neq 2420$ \\ $\rarr$ les coordonnées de $B$ ne vérifient pas l'équation $\quad\rarr\quad$ \textbf{$B$ n'appartient pas à $d$}
		\ei
	\end{correction}
\end{methode}

\begin{rem}
	Pour démontrer que trois points $A$, $B$ et $C$ sont alignés, il suffit de montrer que les coordonnées de $A$ vérifient l'équation de la droite $(BC)$
\end{rem}
\newpage
\subsection{Pente d'une droite}

\begin{prop}{Calcul de la pente}
	Si $A\coordl{x_A}{y_A}$ et $B\coordl{x_B}{y_B}$ sont deux points distincts d'une droite avec $x_A \neq x_B$, la pente est :
	\[
		\boxed{m = \frac{y_B - y_A}{x_B - x_A}}
	\]
\end{prop}

\begin{methode}{Déterminer une équation réduite depuis deux points}
	Soit $A\coordl{4}{-1}$ et $B\coordl{3}{5}$ deux points d'une droite $d$

	Déterminer une équation de $d$

	\begin{correction}
		L'équation réduite est de la forme $y = mx + p$

		\bi
		\item Pente : $$m = \dfrac{5-(-1)}{3-4} = \dfrac{6}{-1} = -6$$
		$\rarr$ L'équation est de la forme $y = -6x + p$
		\item $A$ appartient à $d$ :
		$$
			\begin{array}{rrcl}
				     & y_A & = & -6\times x_A + p \\
				\eqv & -1  & = & -6 \times 4 + p  \\
				\eqv & p   & = & -1 + 24 = 23
			\end{array}$$
		\ei

		L'équation réduite de $d$ est $y = -6x + 23$
	\end{correction}
\end{methode}

\begin{prop}{Droites parallèles et équations réduites}
	Deux droites d'équations $y = mx + p$ et $y = m'x + p'$ sont parallèles si et seulement si $m = m'$
\end{prop}
\begin{methode}{Déterminer la position relative de deux droites}
	Dans chaque cas, déterminer la position relative des deux droites :
	\begin{enumerate}
		\item $d_1 : y = -2x - 5$ et $d_2 : y = -2x + 4$
		\item $d_3 : y = 2x + 1$ et $d_4 : y = -3x + 8$
		\item $d_5 : y = -x + 7$ et $d_6 : y = 3$
		\item $d_7 : x = 1$ et $d_8 : x = -8$
	\end{enumerate}

	\begin{correction}
		\begin{enumerate}
			\item $d_1 \parallel d_2$ car même pente $-2$
			\item $d_3$ et $d_4$ sont \textbf{sécantes} car pentes différentes ($2 \neq -3$).
			\item $d_5$ et $d_6$ sont \textbf{sécantes} car pentes différentes ($-1 \neq 0$).
			\item $d_7 \parallel d_8$ car toutes deux parallèles à l'axe des ordonnées.
		\end{enumerate}
	\end{correction}
\end{methode}
\newpage
% ══════════════════════════════════════════════
\section{Projeté orthogonal d'un point sur une droite}
% ══════════════════════════════════════════════

\begin{defin}{Projeté orthogonal}
	\begin{wrapfigure}{r}{4.5cm}
		\vspace*{-5mm}
		\begin{tikzpicture}[scale=0.9]
			% droite d
			\draw[coulAccent, thick] (0,0) -- (4.5,0) node[right] {$d$};
			% point M
			\coordinate (M) at (2.5,1.8);
			\coordinate (H) at (2.5,0);
			\filldraw[coulAlerte] (M) circle (0.05) node[above] {$M$};
			\filldraw[coulPrimaire] (H) circle (0.05) node[below] {$H$};
			% perpendiculaire
			\draw[dashed, coulBordure] (M) -- (H);
			% angle droit
			\draw (2.5,0.18) -- (2.68,0.18) -- (2.68,0);
		\end{tikzpicture}
	\end{wrapfigure}
	Soit une droite $d$ et un point $M$.

	Le \textbf{projeté orthogonal} du point $M$ sur la droite $d$ est le point d'intersection $H$ de $d$ avec la perpendiculaire à $d$ passant par $M$.
	\vspace*{1.2cm}
\end{defin}

\begin{prop}{Le projeté est le point le plus proche}
	Le projeté orthogonal $H$ du point $M$ sur la droite $d$ est le point de $d$ le plus proche de $M$
\end{prop}

\begin{demo}{au programme}
	Supposons qu'il existe un point $K$ de $d$ tel que $KM \le HM$

	\begin{center}
		\begin{tikzpicture}[scale=1.1]
			\draw[coulAccent, thick] (-0.5,0) -- (4.5,0) node[right] {$d$};
			\coordinate (M) at (2.5,1.8);
			\coordinate (H) at (2.5,0);
			\coordinate (K) at (1.2,0);
			\filldraw[coulAlerte] (M) circle (0.05) node[above] {$M$};
			\filldraw[coulPrimaire] (H) circle (0.05) node[below right] {$H$};
			\filldraw[coulVert] (K) circle (0.05) node[below] {$K$};
			\draw[dashed, coulBordure] (M) -- (H);
			\draw[coulVert] (M) -- (K);
			\draw[coulAlerte] (M) -- (H);
			\draw (2.5,0.18) -- (2.68,0.18) -- (2.68,0);
		\end{tikzpicture}
	\end{center}

	\ms

	$KM \le HM \eqv KM^2 \le HM^2\qquad$ (car $KM\ge 0$ et $HM\ge 0$)

	\ms

	Or le triangle $KHM$ est rectangle en $H$ (par définition du projeté), donc par Pythagore :

	\[
		HM^2 + HK^2 = KM^2
	\]

	Donc $HM^2 + HK^2 \le HM^2$, soit $HK^2 \le 0$

	Ceci est impossible sauf si $HK = 0$, c'est-à-dire $K = H$

	Donc $H$ est bien le point de $d$ le plus proche de $M$
\end{demo}

\ms

\begin{demo}{$\cos^2\alpha + \sin^2\alpha = 1$}
	Soit $d$ une droite, $P \in d$ et $M \notin d$

	On note $H$ le projeté orthogonal de $M$ sur $d$ et $\alpha = \widehat{MPH}$

	\begin{center}
		\begin{tikzpicture}[scale=1.1]
			\draw[coulAccent, thick] (-0.5,0) -- (4,0) node[right] {$d$};
			\coordinate (P) at (0.5,0);
			\coordinate (H) at (3,0);
			\coordinate (M) at (3,2);
			\filldraw[coulPrimaire] (P) circle (0.05) node[below] {$P$};
			\filldraw[coulPrimaire] (H) circle (0.05) node[below] {$H$};
			\filldraw[coulAlerte] (M) circle (0.05) node[right] {$M$};
			\draw[coulPrimaire, thick] (P) -- (M) node[midway, above left] {$PM$};
			\draw[dashed, coulBordure] (M) -- (H)
			node[midway, right] {$HM$};
			\draw[coulAccent] (P) -- (H) node[midway, below] {$PH$};
			\draw (3,0.2) -- (2.8,0.2) -- (2.8,0);
			% angle alpha
			\draw[coulOrange] (0.9,0) arc (0:33:0.4)
			node[right, font=\small] {$\alpha$};
		\end{tikzpicture}
	\end{center}

	Le triangle $PHM$ est rectangle en $H$, donc :
	\[
		\cos\alpha = \frac{PH}{PM} \Rarr PH = PM\times\cos\alpha, \qquad
		\sin\alpha = \frac{HM}{PM} \Rarr HM = PM\times\sin\alpha.
	\]
	Par le théorème de Pythagore :
	\[
		\begin{array}{rccccc}
			      & PH^2                       & + & HM^2                       & = & PM^2 \\
			\Rarr & (PM\times\cos\alpha)^2     & + & (PM\times\sin\alpha)^2     & = & PM^2 \\
			\Rarr & (PM)^2\times(\cos\alpha)^2 & + & (PM)^2\times(\sin\alpha)^2 & = & PM^2 \\
			\Rarr & \cos^2\alpha               & + & \sin^2\alpha               & = & 1
		\end{array}
	\]
\end{demo}

\end{document}
